\documentclass[11pt,a4paper]{article}
\usepackage[T1]{fontenc}
\usepackage[utf8]{inputenc}
\usepackage{amsmath,amsthm,mathtools}
\usepackage{newtxtext,newtxmath}
\usepackage[margin=25mm,headheight=14pt]{geometry}
\usepackage{microtype,booktabs,array,longtable}
\usepackage{fancyhdr,xurl}
\usepackage[hidelinks,breaklinks]{hyperref}
\usepackage{bookmark}
\hypersetup{pdftitle={Poisson Extremes and Gaussian Bridges under Small-Sum Conditioning},
 pdfauthor={Research manuscript prepared for Vladimir Reshetnikov},
 pdfsubject={Polynomial and lacunary uniform series; extremal localization; logarithmic expansions}}
\pagestyle{fancy}\fancyhf{}
\fancyhead[L]{\small Poisson extremes and Gaussian bridges}
\fancyhead[R]{\small ProveIt research extension}
\fancyfoot[C]{\thepage}
\setlength{\parskip}{3pt}
\setlength{\parindent}{1.2em}
\setlength{\emergencystretch}{2em}
\setcounter{tocdepth}{1}
\numberwithin{equation}{section}
\newtheorem{theorem}{Theorem}[section]
\newtheorem{lemma}[theorem]{Lemma}
\newtheorem{proposition}[theorem]{Proposition}
\newtheorem{corollary}[theorem]{Corollary}
\theoremstyle{definition}
\newtheorem{definition}[theorem]{Definition}
\theoremstyle{remark}
\newtheorem{remark}[theorem]{Remark}
\newcommand{\Pp}{\mathbb P}
\newcommand{\E}{\mathbb E}
\newcommand{\R}{\mathbb R}
\newcommand{\N}{\mathbb N}
\newcommand{\ind}{\mathbf 1}
\newcommand{\Var}{\operatorname{Var}}
\newcommand{\Law}{\mathcal L}
\newcommand{\Exp}{\operatorname{Exp}}
\newcommand{\TV}{d_{\mathrm{TV}}}
\newcommand{\supp}{\operatorname{supp}}
\newcommand{\Poi}{\operatorname{PPP}}
\newcommand{\Ga}{\Gamma}
\newcommand{\ph}{\varphi}
\newcommand{\eps}{\varepsilon}
\newcommand{\rise}[2]{(#1)_{#2}}
\newcommand{\repo}{\texttt{ProveIt}}
\newcommand{\commit}{\texttt{a4268e78ebd0f6bf71ba609c07f4b3415748f532}}

\begin{document}
\hypersetup{pageanchor=false}
\begin{titlepage}
\vspace*{7mm}
{\large\sffamily PROVEIT RESEARCH EXTENSION\par}
\vspace{13mm}
{\Huge\bfseries Poisson Extremes and\par Gaussian Bridges under\par Small-Sum Conditioning\par}
\vspace{8mm}
{\Large Localization, incomplete-gamma laws, and\par all-order logarithmic corrections\par}
\vspace{9mm}
{\normalsize Prepared for Vladimir Reshetnikov\par
30 September 2026\par}
\vspace{10mm}
\begin{abstract}
For a positive summable weight sequence, let $S=\sum_j w_jU_j$, with
independent uniform $U_j$ on $[0,1]$, and condition on $S\le x$ as
$x\downarrow0$. We investigate the locations and heights of the largest
summands, jointly with the fluctuation process of the whole configuration.
For $w_j=j^{-p}$, $p>1$, write $N=t^{1/p}$ for the effective dimension at the
mean-matching exponential tilt $t$, put $\beta=1/p$, and define
$b+\beta\log b=\log N$. We prove, uniformly for bounded $z$,
\[
 \log\Pp\!\left(\max_j tj^{-p}U_j\le b+z\mid S\le x\right)
 =-N\Gamma(1-\beta,b+z)+O_p(N^{-1/2}).
\]
Consequently every fixed-order logarithmic correction is explicit. The
extremal point process, with locations divided by $N/b^\beta$, tends to a
Poisson process of intensity $\ind_{[0,1]}(u)\,du\,e^{-z}dz$. It is jointly
independent in the limit of the variance-clock Brownian bridge and the
exponential distance from the conditioning boundary.

For every uniformly lacunary weight sequence, we also prove the full
marked Poisson limit and its joint independence from the bridge and the
boundary distance. This settles the qualitative point-process and
independence portions of Question~7 in the repository's lacunary-conditioning
report; that question's quantitative-rate requirement remains open.
The polynomial extreme-location scale and the all-order maximum law are
additional extensions. Proofs, reproducible symbolic checks, and numerical
Fourier-inversion diagnostics are supplied. No Lean certification,
external refereeing, or exhaustive worldwide priority determination is claimed.
\end{abstract}
\vfill
\noindent\small
\textbf{Inspected repository snapshot.}\quad\commit.\\
\textbf{Research status.}\quad Complete manuscript proofs of the displayed
results; candidate contributions relative to the inspected sources.
\end{titlepage}
\hypersetup{pageanchor=true}
\pagenumbering{roman}
\tableofcontents
\clearpage
\pagenumbering{arabic}

\section{Research target and the main results}
\label{sec:intro}

\subsection{What the repository already establishes}
The Fabius--Rvachev part of \repo\ develops the probability distribution
associated with $\sum_{j\ge1}2^{-j}U_j$. This probabilistic interpretation
is classical; it appears, with the corresponding shift of normalization,
in Arias de Reyna's Theorem~3 \cite{arias}. Our polynomially weighted
series is a related model, not another normalization of the Fabius function.

Two inspected repository manuscripts are particularly relevant. The
lacunary report \cite{repolac} studies full-configuration approximations,
variance-fraction effects, a Brownian bridge, an exponential boundary
slack, and a Gumbel maximum. Its Question~7 asks for the entire extremal
point process and for joint independence from the bridge, with explicit
rates. The report expressly does not claim that independence. A second
manuscript \cite{repounif} already gives a general variance-fraction
conditioning theory for arbitrary positive summable uniform weights,
including polynomial variance clocks. Thus extending that variance
criterion to $j^{-p}$ is \emph{not} a new result of the present article.

The contributions developed here concern a more refined question:
\emph{where do the extremes occur, how accurately can their distribution
be calculated, and are they independent of the macroscopic bridge?}
Theorem~\ref{thm:lacunary} answers the qualitative parts of the repository's
Question~7 throughout its uniformly lacunary class.
Theorems~\ref{thm:gamma} and \ref{thm:jointpoly} supply the distinct
polynomial-weight theory, including a quantitative maximum law strong
enough to imply every fixed-order inverse-logarithmic correction.

Exponential tilting and increasing-block conditional approximation are
classical, not originality claims: see Diaconis--Freedman
\cite{df} and Dembo--Zeitouni \cite{dz}. Poisson limits for independent
rare events and Gaussian bridge conditioning are likewise standard
mechanisms. We give their needed arguments explicitly. The proposed
addition is the theorem package for the specified infinite, increasingly
tilted arrays, the exact incomplete-gamma reduction, and the joint limits.
The literature review was targeted, rather than exhaustive. None of these
statements is identified here with the resolution of a named published
conjecture outside the repository.

\subsection{Notation and the conditioning convention}
All uniforms are independent under the original probability measure $P$.
For a positive summable sequence $(w_j)$, put
\begin{equation}
 S=\sum_{j\ge1}w_jU_j,\qquad
 L(t)=\E_Pe^{-tS},\qquad
 \frac{dQ_t}{dP}=\frac{e^{-tS}}{L(t)}. \label{eq:tilt}
\end{equation}
Choose $t>0$ so that $x=\E_{Q_t}S$, and write
\begin{equation}
 a_j=tw_j,\quad Y_j=a_jU_j,\quad T=\sum_jY_j,\quad
 \mu=\E_{Q_t}T=tx,\quad V=\Var_{Q_t}T,\quad
 P_x=P(\,\cdot\mid S\le x). \label{eq:notation}
\end{equation}
The nonnegative slack under $P_x$ is
$H=\mu-T=t(x-S)$. All conditioning in the main results is on the
positive-probability inequality $S\le x$, not on a null event $S=x$.
The parameter $t$ tends to infinity as $x$ tends to zero.

Let $m_j=\E_{Q_t}Y_j$ and $v_j=\Var_{Q_t}Y_j$. Define the exact variance
clock and its polygonal process by
\begin{equation}
 r_0=0,\qquad r_k=V^{-1}\sum_{j\le k}v_j,\qquad
 \mathcal B_t(r_k)=V^{-1/2}\sum_{j\le k}(Y_j-m_j). \label{eq:clock}
\end{equation}
Linearly interpolate between consecutive $r_k$, and set
$\mathcal B_t(1)=(T-\mu)/\sqrt V$. This gives a continuous function on
$[0,1]$: both $\sum Y_j$ and $\sum m_j$ are finite. A standard Brownian
bridge is denoted $B^\circ$; its covariance is
$\min(r,s)-rs$. The standard normal density is $\ph$.

Point processes below are random locally finite counting measures on
$[0,\infty)\times\R$, with the vague topology. A Poisson point process
with intensity $\ind_{[0,1]}(u)\,du\,e^{-z}dz$ is denoted $\Xi$.
It has finitely many points above any fixed height and arbitrarily many
points towards height $-\infty$. Its boundary lines at $u=0,1$ have zero
intensity. Weak convergence of processes refers to the uniform topology
on $C([0,1])$; convergence of the triples uses the product topology.

\subsection{Polynomial weights: quantitative extremes}
For the rest of this subsection take $w_j=j^{-p}$, where $p>1$ is fixed,
and define
\begin{equation}
 \beta=\frac1p,\qquad N=t^\beta,\qquad
 b+\beta\log b=\log N,\qquad M=Nb^{-\beta}=e^b.
 \label{eq:scales}
\end{equation}
For $N>1$ the equation for $b$ has a unique positive solution. Neither
$N$ nor $M$ needs to be an integer. Also
\begin{equation}
 b=\log N-\beta\log\log N+o(1),\qquad
 M\sim\frac{N}{(\log N)^\beta}. \label{eq:basy}
\end{equation}
We use the upper incomplete gamma function
$\Gamma(s,y)=\int_y^\infty e^{-v}v^{s-1}\,dv$.

\begin{theorem}[Incomplete-gamma maximum law]
\label{thm:gamma}
For every fixed $p>1$ and $Z<\infty$, as $N\to\infty$,
\begin{equation}
 \sup_{|z|\le Z}\left|
 \log P_x\!\left(\max_{j\ge1}Y_j\le b+z\right)
       +N\Gamma(1-\beta,b+z)\right|=O_{p,Z}(N^{-1/2}).
 \label{eq:gamma-main}
\end{equation}
More precisely, the proof gives an error
$O_{p,Z}(N^{-1/2}+b^2/N+e^{-b})$; the last two terms are
$o(N^{-1/2})$. For every fixed integer $K\ge0$,
\begin{align}
 \log P_x\!\left(\max_jY_j\le b+z\right)
 &=-e^{-z}\sum_{r=0}^K
   \frac{(-1)^r\rise{\beta}{r}}{b^r}
        \sum_{k=0}^r\frac{z^k}{k!}
   +O_{p,Z,K}(b^{-K-1}+N^{-1/2}), \label{eq:allorder}
\end{align}
uniformly for $|z|\le Z$, where $\rise\beta0=1$ and
$\rise\beta r=\beta(\beta+1)\cdots(\beta+r-1)$.
\end{theorem}
In particular, the first three terms in the brackets are
\begin{equation}
 1-\frac{\beta(z+1)}b
 +\frac{\beta(\beta+1)(z^2/2+z+1)}{b^2}.
 \label{eq:twocorrections}
\end{equation}
The logarithmic expansion is asymptotic order by order. It is not
claimed to converge, to have a uniform remainder for growing $K$, or to
remain valid when $z$ is unbounded. The theorem implies the Gumbel law
$P_x(\max_jY_j-b\le z)\to\exp(-e^{-z})$.

\begin{theorem}[Locations, bridge, and slack for polynomial weights]
\label{thm:jointpoly}
Let
\begin{equation}
 \Xi_t=\sum_{j\ge1}\delta_{(j/M,\,Y_j-b)}.
 \label{eq:polypp}
\end{equation}
Under $P_x$,
\begin{equation}
  (\mathcal B_t,\Xi_t,H)
       \ \Rightarrow\ (B^\circ,\Xi,E),\qquad E\sim\Exp(1),
 \label{eq:joint-main}
\end{equation}
and the three limiting objects are mutually independent.
If $J_t$ is the almost surely unique index attaining $\max_jY_j$, then
\begin{equation}
 \left(\frac{J_t}{M},\max_jY_j-b\right)
 \Rightarrow (U,G), \label{eq:argmax}
\end{equation}
where $U$ is uniform on $[0,1]$, $P(G\le z)=e^{-e^{-z}}$, and $U,G$
are independent. This pair is also independent of the limiting bridge
and slack.
\end{theorem}
The localization is genuinely submacroscopic: $M/N\to0$, whereas
$J_t/M$ has a nondegenerate limit. A variance clock alone does not reveal
this scale. The upper caps $a_j=(N/j)^p$ force it.

\subsection{Lacunary weights: the repository's qualitative question}
Suppose instead that
\begin{equation}
 w_j>0,\qquad w_{j+1}\le q w_j\quad(j\ge1),\qquad 0<q<1,
 \qquad n=\#\{j:tw_j\ge1\}. \label{eq:lacunarity}
\end{equation}
Set $\Xi_t^{\rm lac}=\sum_j\delta_{(j/n,Y_j-\log n)}$.

\begin{theorem}[Marked extremes for every uniformly lacunary sequence]
\label{thm:lacunary}
Under \eqref{eq:lacunarity}, as $x\downarrow0$,
\begin{equation}
 (\mathcal B_t,\Xi_t^{\rm lac},H)
       \Rightarrow(B^\circ,\Xi,E)
 \label{eq:lacjoint}
\end{equation}
under $P_x$, with the same intensity and mutual independence as in
Theorem~\ref{thm:jointpoly}. In particular the limiting location of the
maximum, on the scale $n$, is uniform and independent of its Gumbel
height, the bridge, and the exponential slack.
\end{theorem}
This includes $w_j=2^{-j}$ and arbitrary supergeometric gaps compatible
with \eqref{eq:lacunarity}. It proves the full point-process convergence
and joint independence asked for in Question~7 of \cite{repolac}.
No quantitative rate in a metric on the joint path--point-process space
is supplied; that portion of the question remains a research problem.

\section{Exponential tilting and the effective dimension}
\label{sec:tilt}

\subsection{Existence, independence, and exact densities}
The sum $S$ lies in $[0,\sum w_j]$. Every sufficiently small $x>0$ has
$P(S\le x)>0$: first make the deterministic tail less than $x/2$, and
then restrict finitely many uniforms so their weighted sum is below
$x/2$. Differentiation under the bounded integral gives
\[
 -\frac{d}{dt}\log L(t)=\E_{Q_t}S,\qquad
 \frac{d}{dt}\E_{Q_t}S=-\Var_{Q_t}S<0.
\]
The tilted mean is continuous, starts at $\frac12\sum w_j$, and tends
to zero. The last assertion follows by dominated convergence for each
coordinate and the deterministic summable tail bound. Hence the
mean-matching $t$ exists and is unique for
$0<x<\frac12\sum w_j$.

Finite-dimensional factorization and dominated convergence show that
under $Q_t$ the $Y_j$ are independent and have densities
\begin{equation}
 p_a(y)=\frac{e^{-y}}{1-e^{-a}}\ind_{(0,a)}(y),\qquad a=a_j.
 \label{eq:pa}
\end{equation}
Integration or differentiation of the normalizing factor gives
\begin{equation}
 m(a)=1-\frac{a}{e^a-1},\qquad
 v(a)=1-\frac{a^2e^a}{(e^a-1)^2},\qquad
 v(a)=m(a)-a m'(a). \label{eq:mv}
\end{equation}
In particular $0<v(a)\le1$. Near zero,
\[
 m(a)=a/2-a^2/12+O(a^4),\qquad v(a)=a^2/12+O(a^4),
\]
and as $a\to\infty$ the laws in \eqref{eq:pa} tend to $\Exp(1)$.

\begin{lemma}[Uniform moments]
\label{lem:moments}
For every integer $r\ge2$ there are positive constants depending only
on $r$ such that, for $Y\sim p_a$ and $X=Y-m(a)$,
\begin{equation}
 c\min(a^2,1)\le v(a)\le C\min(a^2,1),\quad
 \E|X|^r\le C_r\min(a^r,1),\quad
 \E X^4\le C v(a)^2. \label{eq:moments}
\end{equation}
\end{lemma}
\begin{proof}
For $a\le1$, $|X|\le a$, and $v(a)/a^2$ extends continuously to
$1/12$ at zero and stays positive. For $a\ge1$, the density is bounded
above by $(1-e^{-1})^{-1}e^{-y}$, while $m(a)\le1$, giving all uniform
centered moment bounds. The variance has a positive minimum on
$[1,\infty)$, since it is continuous, positive, and tends to one.
Combining these facts proves the final bound as well.
\end{proof}

\subsection{Polynomial constants and the physical endpoint}
\begin{proposition}[The polynomial variance clock]
\label{prop:constants}
For $a_j=(N/j)^p$, $p>1$, define
\begin{align}
 A_p&=\int_0^\infty m(u^{-p})\,du,&
 B_p&=\int_0^\infty v(u^{-p})\,du, \label{eq:AB}\\
 C_p&=\int_0^\infty\log\frac{u^{-p}}{1-e^{-u^{-p}}}\,du.
 \label{eq:C}
\end{align}
These are finite and positive, and
\begin{gather}
 \mu\sim A_pN,\qquad V\sim B_pN,\qquad -\log L(t)\sim C_pN,
 \label{eq:dimension}\\
 A_p=\beta C_p,\qquad B_p=(1-\beta)A_p,\qquad
 C_p=-\Gamma(1-\beta)\zeta(1-\beta). \label{eq:constants}
\end{gather}
For every $r\ge2$, $\sum_j\E_{Q_t}|Y_j-m_j|^r=O_{p,r}(N)$.
The index clock converges to
\begin{equation}
 F_p(u)=\frac1{B_p}\int_0^u v(s^{-p})\,ds,\quad
 r_{\lfloor Nu\rfloor}\to F_p(u),\qquad 0\le u<\infty.
 \label{eq:indexclock}
\end{equation}
Moreover
\begin{equation}
 x\sim A_p t^{\beta-1},\quad
 t\sim(A_p/x)^{1/(1-\beta)},\quad
 N\sim(A_p/x)^{1/(p-1)}. \label{eq:xscale}
\end{equation}
\end{proposition}
\begin{proof}
Riemann sums apply after separately controlling $j<\delta N$ and
$j>RN$. For $m$ and $v$, the small-index summands are at most one; the
large-index tails are bounded respectively by constants times
$(N/j)^p$ and $(N/j)^{2p}$. For the logarithmic integrand the behavior
at $u=0$ is $-p\log u+o(1)$ and at infinity it is
$u^{-p}/2+O(u^{-2p})$. The endpoint contributions are integrable and
uniformly negligible as $\delta\downarrow0$ and $R\uparrow\infty$.
This proves the three asymptotics and the clock limit. Lemma~\ref{lem:moments}
proves the moment sum by the same head--tail split.

Put $g(a)=\log[a/(1-e^{-a})]$. Then $g'(a)=m(a)/a$.
The substitution $a=u^{-p}$ and integration by parts give
\[
 C_p=\beta\int_0^\infty g(a)a^{-\beta-1}\,da
     =\int_0^\infty m(a)a^{-\beta-1}\,da=A_p/\beta.
\]
All boundary terms vanish: $g(a)=O(a)$ at zero and $g(a)=O(\log a)$
at infinity. Similarly, using \eqref{eq:mv},
\[
 B_p=\beta\int_0^\infty[m(a)-am'(a)]a^{-\beta-1}\,da
     =(1-\beta)A_p.
\]
For completeness, the regularized Euler integral for $0<s<1$ is
\begin{equation}
 \Gamma(s)\zeta(s)=\int_0^\infty
 \left(\frac1{e^a-1}-\frac1a\right)a^{s-1}\,da.
 \label{eq:zeta-reg}
\end{equation}
It follows from the usual Euler integral for $s>1$ by splitting at one,
subtracting $1/a$ on $(0,1)$, and analytically continuing the explicit
term $1/(s-1)$; on $0<s<1$ that term equals the integral of $-1/a$
over $(1,\infty)$. The subtracted integral is absolutely convergent.
Taking $s=1-\beta$ and using
$m(a)=-a[(e^a-1)^{-1}-a^{-1}]$ proves the last identity in
\eqref{eq:constants}. Standard gamma and zeta integral conventions are
as in \cite{dlmf}. Finally $x=\mu/t$ proves \eqref{eq:xscale}.
\end{proof}
For $p=2$ the constants are
\[
 C_2=2.588410972826785\ldots,\quad
 A_2=1.294205486413393\ldots,\quad
 B_2=0.647102743206696\ldots.
\]
These scalar facts are background for the new extremal results, not a
claim to supersede the polynomial variance analysis in \cite{repounif}.

\subsection{The lacunary dimension}
\begin{proposition}
\label{prop:lacdimension}
Under \eqref{eq:lacunarity}, $n\to\infty$ and
\begin{equation}
 \mu=n+O_q(1),\qquad V=n+O_q(1),\qquad
 \sum_j\E|Y_j-m_j|^r=O_{q,r}(n),\quad r\ge2.
 \label{eq:lacdimension}
\end{equation}
Uniformly in $k\ge0$,
\begin{equation}
 \sum_{j\le k}v_j=\min(k,n)+O_q(1).
 \label{eq:lacclock}
\end{equation}
\end{proposition}
\begin{proof}
For $j\le n$, $a_j\ge q^{-(n-j)}$. For $k\ge1$,
$a_{n+k}<q^{k-1}$. The deficits $1-m(a)$ and $1-v(a)$ on the active
part are summable along the first bound, because they decay exponentially
as $a\to\infty$; they can also be bounded uniformly for $a\ge1$.
The tails $m(a)=O(a)$ and $v(a)=O(a^2)$ are summable along the second
bound. The same argument and Lemma~\ref{lem:moments} give the moment
estimate and the uniform clock assertion. Every fixed coordinate
becomes active as $t\to\infty$, proving $n\to\infty$.
\end{proof}

\section{A quantitative local limit theorem and conditioning identities}
\label{sec:local}

We give the analytic estimates used below instead of relying on an
unproved transfer of an independent maximum law through a rare event.
Write $R=N$ in the polynomial case and $R=n$ in the lacunary case.
Constants may depend on the fixed parameter $p$ or $q$.

\begin{lemma}[A smoothing block]
\label{lem:smoothing}
For either array, every coordinate set $I$ with
$V_I:=\sum_{j\in I}v_j\ge\eps R$ contains at least $c_\eps R$
coordinates whose caps are at least $a_\eps>0$. The same assertion
holds, uniformly in $y\ge2$, for the polynomial array with caps
$c_j=\min(a_j,y)$ and with its own variances.
\end{lemma}
\begin{proof}
In the polynomial case the first $\delta N$ coordinates contribute
at most $\delta N+1$ variance. The part beyond $AN$ contributes at
most $C_pNA^{1-2p}$, by Lemma~\ref{lem:moments}; here $A$ is a
fixed large constant. Choose $\delta$
small and $A$ large so their combined contribution is less than
$\eps N/2$. At least $\eps N/2$ variance remains in
$[\delta N,AN]$, so at least $\eps N/2$ coordinates remain there,
since each variance is at most one. Their original caps are at least
$A^{-p}$. Capping at $y\ge2$ preserves a positive lower bound on
those caps and cannot increase the two discarded variance bounds.

In the lacunary case the entire inactive tail has $O_q(1)$ variance.
Thus $I$ contains at least $\eps n/2$ active coordinates for large $n$,
and all their caps are at least one.
\end{proof}

\begin{lemma}[Uniform density local limit theorem]
\label{lem:llt}
Let $I$ be any coordinate set with $V_I\ge\eps R$. The sum
$T_I=\sum_{j\in I}Y_j$ has a continuous density $f_I$, and
\begin{equation}
 \sup_{z\in\R}
 \left|\sqrt{V_I}\,f_I(\mu_I+\sqrt{V_I}z)-\ph(z)\right|
 \le C_\eps R^{-1/2},\qquad \mu_I=\E T_I.
 \label{eq:llt}
\end{equation}
This is also uniform for the polynomial capped arrays of
Lemma~\ref{lem:smoothing}. In particular
$\sup f_I\le C_\eps R^{-1/2}$.
\end{lemma}
\begin{proof}
For a centered variable with density \eqref{eq:pa}, write $\chi_a(s)$
for its characteristic function. Apart from its unit-modulus centering
factor, it is
\begin{equation}
 \frac{1-e^{-(1-is)a}}{(1-e^{-a})(1-is)}. \label{eq:cf}
\end{equation}
Uniformly for $a\ge a_0>0$, this has modulus at most $C_{a_0}/|s|$
for $|s|\ge1$. On any compact frequency interval
$\eta\le|s|\le A$ it has modulus at most some $\rho<1$,
uniformly for $a\ge a_0$. To see strict uniformity, adjoin $a=\infty$:
the expression converges there to $(1-is)^{-1}$, while at every finite
positive $a$ a variable with a positive interval density has
characteristic-function modulus strictly below one away from zero.
Compactness then applies.

The moment bounds also imply, uniformly for sufficiently small $|s|$,
\begin{equation}
 \log\chi_a(s)=-\tfrac12v(a)s^2+
 O\!\left(\min(a^3,1)|s|^3\right). \label{eq:cumulant}
\end{equation}
For $a\le1$ this follows by Taylor expansion of a bounded centered
variable, whose variance is comparable to $a^2$; for $a\ge1$ the
uniform exponential moment bound supplies the same expansion on a
fixed neighborhood of zero. The logarithm is taken on its continuous
branch equal to zero at $s=0$.

Let $\Psi_I(u)$ be the characteristic function of
$(T_I-\mu_I)/\sqrt{V_I}$. The sum of the error coefficients in
\eqref{eq:cumulant} is $O(R)$, and $V_I\asymp_\eps R$.
For $|u|\le R^{1/12}$,
\[
 \Psi_I(u)=e^{-u^2/2}
 \left[1+O_\eps(R^{-1/2}|u|^3)
                e^{O_\eps(R^{-1/2}|u|^3)}\right].
\]
Its integrated difference from $e^{-u^2/2}$ on that interval is
$O_\eps(R^{-1/2})$. Between $R^{1/12}$ and
$\eta\sqrt{V_I}$, the small-frequency modulus bound from
\eqref{eq:cumulant}, with $\eta$ chosen sufficiently small, gives
$|\Psi_I(u)|\le e^{-c_\eps u^2}$. Between
$\eta\sqrt{V_I}$ and $A\sqrt{V_I}$ the smoothing block gives
$|\Psi_I(u)|\le\rho^{c_\eps R}$. Above $A\sqrt{V_I}$ it gives
\[
 |\Psi_I(u)|\le
 \left(\frac{C\sqrt{V_I}}{|u|}\right)^{c_\eps R}.
\]
Choose $A>2C$. Integrating these three bounds, and the analogous
Gaussian tails, contributes an exponentially small quantity.
Thus $\int_\R|\Psi_I(u)-e^{-u^2/2}|du=O_\eps(R^{-1/2})$.

For large $R$ the smoothing block has more than two coordinates, so
its characteristic function is integrable. The remaining independent
coordinates have characteristic-function modulus at most one. The
same argument therefore applies to infinite $I$, by convergence of the
sums and dominated Fourier inversion. The inversion inequality now
proves \eqref{eq:llt} and the density bound.
\end{proof}

\subsection{The exact likelihood and the slack}
Let $f$ be the tilted density of $T$, and define
\begin{equation}
 D=\E_{Q_t}\big[e^{T-\mu}\ind_{T\le\mu}\big]
   =\int_0^\infty e^{-h}f(\mu-h)\,dh.
 \label{eq:D}
\end{equation}
Changing measure in \eqref{eq:tilt} gives the exact identities
\begin{equation}
 \frac{dP_x}{dQ_t}=D^{-1}e^{T-\mu}\ind_{T\le\mu},
 \qquad P(S\le x)=L(t)e^\mu D. \label{eq:likelihood}
\end{equation}
It is important that this full-product likelihood is not close to one.
Approximation of selected marginals requires integrating out the
complement.

\begin{proposition}[Normalization and boundary slack]
\label{prop:slack}
In either array,
\begin{equation}
 D=(2\pi V)^{-1/2}\big(1+O(R^{-1/2})\big). \label{eq:saddle}
\end{equation}
Under $P_x$, $H$ tends in total variation to $\Exp(1)$, and
$\E_{P_x}H^r\to r!$ for every fixed nonnegative integer $r$.
In the polynomial case,
\begin{equation}
 -\log P(S\le x)\sim (1-\beta)C_pN. \label{eq:smallball}
\end{equation}
\end{proposition}
\begin{proof}
Apply Lemma~\ref{lem:llt} to $T$ in \eqref{eq:D}. After multiplying
by $\sqrt V$, the integral of the error is $O(R^{-1/2})$ and
\[
 \int_0^\infty e^{-h}\ph(h/\sqrt V)\,dh
 =(2\pi)^{-1/2}\big(1+O(V^{-1})\big).
\]
This proves \eqref{eq:saddle}. The density of $H$ is exactly
$e^{-h}f(\mu-h)/D$, $h\ge0$. It tends pointwise to $e^{-h}$ and
is bounded by $Ce^{-h}$, by the density bound in
Lemma~\ref{lem:llt}. Dominated convergence proves both total variation
and the moments. Finally use \eqref{eq:likelihood},
Proposition~\ref{prop:constants}, and $\log D=O(\log N)$.
\end{proof}

For $I\subseteq\N$ put $Z_I=T_I-\mu_I$ and $J=I^c$. When
$V_J\ge\eps R$, the joint law of $((Y_j)_{j\in I},H)$ has density
relative to $Q_t^I\otimes dh$ equal to
\begin{equation}
 R_I(\mathbf y,h)=
 e^{-h}\frac{f_J(\mu_J-Z_I(\mathbf y)-h)}D\ind_{h\ge0}.
 \label{eq:jointlikelihood}
\end{equation}
This follows by disintegration of the independent sums; it applies to
infinite coordinate sets because their nonnegative sums are measurable
limits. In particular $R_I(\mathbf y,h)\le C_\eps e^{-h}$.

\begin{corollary}[Negligible blocks with independent slack]
\label{cor:smallblock}
If $V_I=o(R)$, then
\begin{equation}
 \TV\!\left(\Law_{P_x}((Y_j)_{j\in I},H),
             Q_t^I\otimes\Exp(1)\right)\longrightarrow0.
 \label{eq:smallblock}
\end{equation}
More generally, if $V_{I^c}\ge\eps R$, the marginal likelihood of
$P_x^I$ with respect to $Q_t^I$ is bounded by $C_\eps$.
\end{corollary}
\begin{proof}
The general bound follows by integrating
\eqref{eq:jointlikelihood} in $h$. In the negligible-block case,
$Z_I/\sqrt R\to0$ in $Q_t$-probability by Chebyshev's inequality,
and $V_{I^c}/V\to1$. Lemma~\ref{lem:llt} and
\eqref{eq:saddle} show that the ratio in
\eqref{eq:jointlikelihood} tends to one in probability for each fixed
$h$. Its uniform bound, followed by integration against $e^{-h}dh$,
gives convergence in $L^1(Q_t^I\otimes\Exp(1))$. Half this $L^1$
distance is total variation.
\end{proof}
This special case of the general variance principle is reproduced for
the proofs below; the broader principle is already in \cite{repounif}.

\section{The incomplete-gamma law for the maximum}
\label{sec:maximum}

The proof separates an exact independent-array calculation from a
quantitative conditioning estimate. Merely knowing that an extreme
prefix has negligible variance would give an $o(1)$ error, insufficient
for every logarithmic correction.

\subsection{An exact intensity integral}
For $y>0$, the tilted exceedance probability at coordinate $j$ is
\begin{equation}
 p_j(y)=\frac{e^{-y}-e^{-a_j}}{1-e^{-a_j}}\ind_{a_j>y}.
 \label{eq:exceedance}
\end{equation}
Only finitely many coordinates can exceed $y$. Put
$q_j(y)=(e^{-y}-e^{-a_j})_+$ and
$\Lambda_N(y)=\sum_jp_j(y)$.

\begin{lemma}[Integral and lattice errors]
\label{lem:gammaint}
For $a_j=(N/j)^p$,
\begin{equation}
 \int_0^\infty\big(e^{-y}-e^{-(N/s)^p}\big)_+\,ds
       =N\Gamma(1-\beta,y). \label{eq:gammaint}
\end{equation}
Furthermore
\begin{equation}
 \left|\sum_{j\ge1}q_j(y)-N\Gamma(1-\beta,y)\right|
       \le e^{-y}. \label{eq:lattice}
\end{equation}
Uniformly for $y=b+z$, $|z|\le Z$,
\begin{equation}
 \log Q_t(\max_jY_j\le y)
       =-N\Gamma(1-\beta,y)+O_Z(e^{-b}).
 \label{eq:tiltedmax}
\end{equation}
\end{lemma}
\begin{proof}
The integrand in \eqref{eq:gammaint} is supported on
$0<s<Ny^{-\beta}$. The substitution $a=(N/s)^p$ gives
\[
 Ny^{-\beta}e^{-y}
 -\beta N\int_y^\infty e^{-a}a^{-\beta-1}\,da.
\]
An integration by parts identifies this with
$N\int_y^\infty e^{-a}a^{-\beta}\,da$. As a function of $s$, the
original integrand is nonincreasing, takes limiting value $e^{-y}$
at zero, and vanishes eventually. The elementary integral comparison
for a nonincreasing nonnegative function therefore gives
\eqref{eq:lattice}.

For contributing coordinates $a_j>y$,
\[
 0\le p_j(y)-q_j(y)
 \le\frac{e^{-y}}{1-e^{-y}}q_j(y).
\]
When $y=b+z$ with $z$ bounded, \eqref{eq:gammaint} and
$Ne^{-b}b^{-\beta}=1$ give
$\sum q_j=O_Z(1)$, and hence
$\Lambda_N(y)=N\Gamma(1-\beta,y)+O_Z(e^{-b})$.
Also $\max_jp_j(y)\le e^{-y}/(1-e^{-y})=O_Z(e^{-b})$.
The independent-product formula and
$|\log(1-u)+u|\le u^2/[2(1-u)]$ show that
\[
 \log Q_t(\max_jY_j\le y)
 =-\Lambda_N(y)+O\!\left(\sum_jp_j(y)^2\right).
\]
The last sum is at most $\max_jp_j\cdot\Lambda_N=O_Z(e^{-b})$.
\end{proof}

\subsection{Removing the small-sum constraint at logarithmic accuracy}
Let $Q_t^{[y]}=Q_t(\,\cdot\mid\max_jY_j\le y)$. This is an
independent product with caps $c_j=\min(a_j,y)$, since the event
restricts each of only finitely many coordinates separately. Denote
its total mean, variance, and density by $\mu_y,V_y,f_y$. Define
\begin{equation}
 D_y=\int_0^\infty e^{-h}f_y(\mu-h)\,dh. \label{eq:Dy}
\end{equation}
Notice the \emph{original} mean $\mu$ in this definition. Then
\begin{equation}
 P_x(\max_jY_j\le y)
    =Q_t(\max_jY_j\le y)\frac{D_y}{D}. \label{eq:capidentity}
\end{equation}

\begin{lemma}[Capping perturbs the normalization negligibly]
\label{lem:capping}
Uniformly for $y=b+z$, $|z|\le Z$,
\begin{equation}
 \delta_y:=\mu-\mu_y=O_Z(y),\qquad
 |V-V_y|=O_Z(y^2), \label{eq:capmoments}
\end{equation}
and
\begin{equation}
 \frac{D_y}{D}=1+O_{p,Z}(N^{-1/2}+y^2/N). \label{eq:capratio}
\end{equation}
\end{lemma}
\begin{proof}
There are $K_y=\lfloor Ny^{-\beta}\rfloor$ possibly changed
coordinates. For $a>y\ge2$,
\[
 0\le m(a)-m(y)
 =\frac{y}{e^y-1}-\frac{a}{e^a-1}\le Cye^{-y},
 \qquad |v(a)-v(y)|\le Cy^2e^{-y}.
\]
For the variance bound use \eqref{eq:mv} and the bound
$a^2e^{-a}\le y^2e^{-y}$ for $a\ge y\ge2$, adjusting the absolute
constant for the denominators. Since $K_ye^{-y}=O_Z(1)$,
\eqref{eq:capmoments} follows. Thus $V_y\sim B_pN$.

Lemma~\ref{lem:llt}, uniformly for the capped array, yields
\[
 D_y=\frac1{\sqrt{2\pi V_y}}
 \int_0^\infty e^{-h}
       \exp\!\left(-\frac{(\delta_y-h)^2}{2V_y}\right)dh
       +O_{p,Z}(N^{-1}).
\]
The integral is $1+O((\delta_y^2+1)/V_y)$, using
$0\le1-e^{-u}\le u$. Compare with
\eqref{eq:saddle} and use \eqref{eq:capmoments} to obtain
\eqref{eq:capratio}.
\end{proof}

\begin{proof}[Proof of Theorem~\ref{thm:gamma}]
Equations \eqref{eq:capidentity}, \eqref{eq:tiltedmax}, and
\eqref{eq:capratio} give
\[
 \log P_x(\max_jY_j\le b+z)
 =-N\Gamma(1-\beta,b+z)
   +O_{p,Z}(e^{-b}+N^{-1/2}+b^2/N).
\]
Because $e^{-b}=b^\beta/N$ and $b\sim\log N$, this proves
\eqref{eq:gamma-main}.

The expansion has a particularly simple direct proof. Write
$a=b+z+s$ in the gamma integral and use the defining equation for $b$:
\begin{equation}
 N\Gamma(1-\beta,b+z)
 =e^{-z}\int_0^\infty e^{-s}
       \left(1+\frac{z+s}{b}\right)^{-\beta}ds.
 \label{eq:gamma-shift}
\end{equation}
Taylor's theorem, for fixed $K$ and bounded $z$, gives an integrated
remainder $O_{\beta,Z,K}(b^{-K-1})$. To check uniformity even for
negative $z$, take $b>2Z$: the Taylor segment stays at least $1/2$
away from the singularity at $-1$, and the remainder is bounded by
$C b^{-K-1}(|z|+s)^{K+1}$, integrable against $e^{-s}$.
The coefficient of $b^{-r}$ is
\[
 \frac{(-1)^r\rise\beta r}{r!}
   \int_0^\infty e^{-s}(z+s)^r ds
 =(-1)^r\rise\beta r\sum_{k=0}^r\frac{z^k}{k!}.
\]
This proves \eqref{eq:allorder}, without an interchange of an infinite
asymptotic series and an integral.
\end{proof}

\begin{remark}[What the error does and does not mean]
The $O(N^{-1/2})$ bound is deliberately conservative. It is adequate
for every fixed inverse-logarithmic order because
$N^{-1/2}=o(b^{-K})$ for every fixed $K$. It is not a claim of a
sharp conditioning correction at the polynomial scale $1/N$.
At moderate $N$ the conditioning error can exceed several logarithmic
terms; the numerical results in Section~\ref{sec:numerics} illustrate
this explicitly.
\end{remark}

\section{Poisson localization for polynomial weights}
\label{sec:poisson}

\begin{proposition}[The tilted marked process]
\label{prop:polypp}
Under $Q_t$, the process $\Xi_t$ in \eqref{eq:polypp} converges to
$\Xi$. The same convergence holds under $P_x$, jointly with an
independent $\Exp(1)$ slack.
\end{proposition}
\begin{proof}
Let $g$ be a nonnegative continuous compactly supported function on
$[0,\infty)\times\R$, with height support in $[L_0,L_1]$. Independence
gives
\[
 \E_{Q_t}e^{-\langle g,\Xi_t\rangle}
   =\prod_j(1-q_{j,t}),\qquad
 q_{j,t}=\E_{Q_t}\big[1-e^{-g(j/M,Y_j-b)}\big].
\]
Only finitely many $q_{j,t}$ are nonzero. On their height support,
$Y_j=b+z$ and $a_j=b(j/M)^{-p}$. Since $e^{-b}=1/M$,
\begin{align}
 \sum_jq_{j,t}
  =\frac1M\sum_j\int_{L_0}^{L_1}
 &(1-e^{-g(j/M,z)})e^{-z}\nonumber\\[-2mm]
 &\times\frac{\ind_{b+z<b(j/M)^{-p}}}
               {1-e^{-b(j/M)^{-p}}}\,dz.
 \label{eq:pp-riemann}
\end{align}
Where the indicator is nonzero, the denominator tends to one
uniformly. For $u<1$ the indicator tends to one; for $u>1$ it tends
to zero. Split the location integral into the two regions away from
$u=1$ and a strip of width $2\delta$ around one. Ordinary Riemann
sums give convergence on the first two regions, and a uniform bounded
integrand controls the strip by $C\delta$. Thus
\[
 \sum_jq_{j,t}\longrightarrow
 \int_0^1\int_\R(1-e^{-g(u,z)})e^{-z}\,dz\,du.
\]
Also $\max_jq_{j,t}=O_g(M^{-1})$, so
$\sum_jq_{j,t}^2\to0$. Taking logarithms proves convergence of
Laplace functionals to that of $\Xi$. The same bounds on compact
sets give local tightness, hence point-process convergence.

For the conditioned assertion, let $I_t=\{1,\ldots,\lceil2M\rceil\}$.
Its tilted variance is at most $2M+1=o(N)$. For each fixed lower
height $L_0$, every coordinate outside $I_t$ has cap at most
$b/2^p<b+L_0$ for sufficiently large $N$. Therefore every point of
$\Xi_t$ above height $L_0$ is determined by $I_t$. Apply
Corollary~\ref{cor:smallblock} and then the tilted convergence to
obtain the joint Poisson and independent-slack limit. Exhausting the
height axis by compact intervals completes the argument.
\end{proof}

\begin{corollary}[All fixed upper order statistics]
\label{cor:orderstats}
Order the contributions as $Y_{(1)}>Y_{(2)}>\cdots$, with indices
$J_{(1)},J_{(2)},\ldots$. For each fixed $r$,
\begin{equation}
 \big((J_{(k)}/M,Y_{(k)}-b)\big)_{k=1}^r
 \Rightarrow \big((U_k,-\log A_k)\big)_{k=1}^r,
 \label{eq:orderstats}
\end{equation}
where $A_k=E_1+\cdots+E_k$, the $E_i$ are independent $\Exp(1)$,
and the $U_k$ are independent uniforms on $[0,1]$, independent of all
$E_i$. In particular
\begin{equation}
 P_x(Y_{(r)}-b\le z)\longrightarrow
 e^{-e^{-z}}\sum_{k=0}^{r-1}\frac{e^{-kz}}{k!}.
 \label{eq:rth}
\end{equation}
\end{corollary}
\begin{proof}
The intensity is a product of unit Lebesgue measure on $[0,1]$ and
$e^{-z}dz$. Under the substitution $v=e^{-z}$, the height intensity
becomes Lebesgue measure on $(0,\infty)$, whose successive Poisson
arrival times are $A_k$. Attaching independent uniform marks gives
the displayed limiting process and hence the limiting order statistics.

To justify passage from vague convergence to top points, there is no
escape in location above any fixed height: all such points lie in
$[0,2]$ for large $N$. Theorem~\ref{thm:gamma} gives tightness of the
highest height. The number above a low level $L_0$ tends to a Poisson
variable with mean $e^{-L_0}$, which exceeds $r$ with probability tending
to one as $L_0\to-\infty$. On a compact rectangle containing the top
$r$ points, the limiting process almost surely has no height ties or
boundary points, so ordering is continuous. The formula
\eqref{eq:rth} is the probability of at most $r-1$ points above $z$.
\end{proof}

\begin{corollary}[No single summand carries the conditioned mass]
\label{cor:nocondensation}
For polynomial weights,
\begin{equation}
 \frac{S}{x}\longrightarrow1,
 \qquad
 \frac{A_pN}{b}\frac{\max_j j^{-p}U_j}{S}\longrightarrow1
 \quad\text{in }P_x\text{-probability}.
 \label{eq:nocondensation}
\end{equation}
In particular the largest mass fraction tends to zero.
\end{corollary}
\begin{proof}
By Proposition~\ref{prop:slack}, $T=\mu-H=\mu+O_{P_x}(1)$,
while $\mu\sim A_pN$. Theorem~\ref{thm:gamma} gives
$\max_jY_j=b+O_{P_x}(1)$. Divide these identities, using $b\to\infty$.
\end{proof}
The extremes are concentrated in a vanishing fraction of the relevant
indices, yet each contributes only order $\log N/N$ of the total mass.
This distinction between location concentration and mass concentration
is useful when interpreting the conditioned configurations.

\section{The Gaussian bridge and joint independence}
\label{sec:bridge}

\subsection{A variance-clock invariance principle}
\begin{lemma}[Tilted motion and conditioned bridge]
\label{lem:bridge}
In both the polynomial and uniformly lacunary cases,
\begin{equation}
 \mathcal B_t\Rightarrow B\quad\text{under }Q_t,
 \qquad
 (\mathcal B_t,H)\Rightarrow(B^\circ,E)
       \quad\text{under }P_x, \label{eq:bridges}
\end{equation}
where $B$ is standard Brownian motion and $B^\circ,E$ are independent.
\end{lemma}
\begin{proof}
Write $X_j=Y_j-m_j$. Since $\max_jv_j\le1$ and $V\asymp R$,
the largest clock interval is $O(R^{-1})$. Lemma~\ref{lem:moments}
gives
\[
 V^{-2}\sum_j\E X_j^4
 \le C V^{-2}\sum_jv_j^2\le C/V\longrightarrow0.
\]
Taylor expansion of characteristic functions therefore proves the
Gaussian finite-dimensional limits for disjoint partial-sum increments.
Their limiting covariance in the exact clock is $\min(r,s)$.
Interpolation at any fixed time changes the covariance by at most the
largest clock interval and changes a random variable by $o_P(1)$.

For tightness, an increment between times $s<r$ is
$V^{-1/2}\sum_jc_jX_j$, with $0\le c_j\le1$ and
$\sum_jc_jv_j=V(r-s)$. Independence and the fourth-moment bound give
\begin{align*}
 \E_{Q_t}|\mathcal B_t(r)-\mathcal B_t(s)|^4
 &\le C V^{-2}\left(\sum_jc_j^2v_j\right)^2\\
 &\le C(r-s)^2.
\end{align*}
This holds also for countably many summands by moment convergence.
The fourth-moment tightness criterion on continuous paths follows,
for example, by Markov's inequality on dyadic increments followed by
the usual summable modulus-of-continuity bound. It proves the first
convergence in \eqref{eq:bridges}.

Fix $0<r_1<\cdots<r_d=c<1$, and take the observed prefix $I$ ending
at the last clock knot before $c$. The clock mesh is negligible and
$V_{I^c}/V\to1-c$. Include $Z_I/\sqrt V$ as the last coordinate in
the observed Gaussian vector. Lemma~\ref{lem:llt},
\eqref{eq:jointlikelihood}, and \eqref{eq:saddle} show that, for
bounded $h$, the joint likelihood is asymptotic to
\begin{equation}
 e^{-h}g_c(Z_I/\sqrt V),\qquad
 g_c(z)=\frac1{\sqrt{1-c}}
          \exp\!\left(-\frac{z^2}{2(1-c)}\right).
 \label{eq:gaussianweight}
\end{equation}
The approximation is uniform in $Z_I/\sqrt V$. The uniform bound
$C_c e^{-h}$ permits integration over all $h\ge0$.
Weighting the Brownian finite-dimensional law by $g_c(B(c))$ gives
the Brownian bridge law up to $c$: indeed $g_c$ is the ratio of the
normal density for completing $B(c)$ to zero in variance $1-c$ to
the unit-variance density at zero. Direct Gaussian conditioning gives
covariance $\min(r,s)-rs$. The independent factor $e^{-h}$ proves
finite-dimensional independence of the slack.

For conditioned tightness on $[0,1-\eps]$, use the bounded prefix
likelihood in Corollary~\ref{cor:smallblock} and the tilted tightness.
To control the terminal interval, take the tail after the last knot
before $1-\eps$, with $0<\eps<1/2$. Its variance is at most
$\eps V+1$ and its complement has a fixed positive fraction of $V$.
Its marginal likelihood is therefore bounded by a constant independent
of $\eps$. Kolmogorov's maximal inequality for reversed tail sums,
first for finite tails and then by passage to the convergent series,
gives
\begin{equation}
 \limsup_{t\to\infty}P_x\left(
 \sup_{r\in[1-\eps,1]}
 |\mathcal B_t(r)-\mathcal B_t(1)|>\eta\right)
 \le C\eps/\eta^2. \label{eq:terminaltight}
\end{equation}
One may replace reverse sums by a terminal sum minus forward partial
sums to obtain the same bound. Interpolation cannot enlarge the
maximum of the absolute endpoint values. Finally
$\mathcal B_t(1)=-H/\sqrt V\to0$ in probability. Letting
$\eps\downarrow0$ proves tightness at the endpoint. The slack is
tight by Proposition~\ref{prop:slack}. The finite-dimensional limits
identify the joint limit uniquely.
\end{proof}

\subsection{Why rare points are independent of Gaussian fluctuations}
The following factorization makes the required independence explicit.
Neither vanishing pairwise correlations nor separate marginal limits
would suffice.

\begin{lemma}[Rare-mark factorization]
\label{lem:factorization}
Consider either tilted array and a marked process centered at $c_t$,
where $c_t=b$ for polynomial weights and $c_t=\log n$ for lacunary
weights. Suppose its compactly supported nonnegative test function
$g$ satisfies
\[
 q_j=\E[1-e^{-g(\text{location}_j,Y_j-c_t)}],\qquad
 \sum_jq_j=O_g(1).
\]
For any bounded deterministic coefficients $\theta_j$,
\begin{align}
 &\E\left[\exp\!\left(i\sum_j\theta_jX_j/\sqrt V\right)
                  e^{-\langle g,\Xi_t\rangle}\right]
 \nonumber\\
 &\hspace{8mm}{}-
 \E\exp\!\left(i\sum_j\theta_jX_j/\sqrt V\right)
       \E e^{-\langle g,\Xi_t\rangle}\longrightarrow0.
 \label{eq:factorization}
\end{align}
The assertion also holds with the marks restricted to any deterministic
coordinate subset.
\end{lemma}
\begin{proof}
Write $\phi_j=\E e^{i\theta_jX_j/\sqrt V}$. The $j$th factor in
the first expectation is
\[
 \phi_j-q_j-e_j,\qquad
 |e_j|\le\frac{C}{\sqrt V}
     \E\big[|X_j|(1-e^{-g(\text{location}_j,Y_j-c_t)})\big].
\]
On the height support of $g$, $Y_j=c_t+O_g(1)$ and $0\le m_j\le1$,
so $\sum_j|e_j|=O_g((c_t+1)/\sqrt V)=o(1)$.
Also $|\phi_j-1|\le C v_j/V\le C/V$, by centering and
$|e^{iu}-1-iu|\le u^2/2$. Replacing the $j$th factor by
$\phi_j(1-q_j)$ has total error at most
\[
 \sum_j|e_j|+\sum_jq_j|\phi_j-1|=o(1).
\]
The telescoping product inequality is applicable since every exact
expectation factor, as well as $\phi_j(1-q_j)$, has modulus at most
one. Independence factors the resulting product into the two
expectations in \eqref{eq:factorization}. Finite products suffice for
the marks; convergence of the centered sums handles an infinite
Gaussian functional.
\end{proof}
Taking $\theta_j$ to be the finite linear combination of partial-sum
indicators corresponding to any chosen clock times proves joint
independence of the tilted Brownian limit and the limiting marked
Poisson process, whenever the latter has been established. Marginal
tightness then gives joint tightness. The argument survives restriction
of the marks to a prefix, which is crucial for conditioning.

\subsection{Proof of the polynomial joint theorem}
\begin{proof}[Proof of Theorem~\ref{thm:jointpoly}]
Fix a compact height window for the marks and clock times
$0<r_1<\cdots<r_d<1$. All marks in that window belong to the first
$\lceil2M\rceil$ coordinates for large $N$. The variance of that
prefix is $o(V)$, so it lies before any fixed positive clock time.
Choose $c\in(r_d,1)$ and observe the entire coordinate prefix up to
clock $c$, including its centered total $Z_I/\sqrt V$.

Under $Q_t$, Proposition~\ref{prop:polypp}, Lemma~\ref{lem:bridge},
and Lemma~\ref{lem:factorization} give the joint limit of the observed
Gaussian vector and the marks as independent Brownian motion and
$\Xi$. Apply the exact conditional likelihood
\eqref{eq:jointlikelihood}. As in \eqref{eq:gaussianweight}, its
limit is $e^{-h}g_c(B(c))$. This weight depends on the Gaussian
component alone and on the separate slack variable, not on $\Xi$.
It therefore changes Brownian motion into a Brownian bridge while
leaving the Poisson process independent and unchanged, and produces
an independent $\Exp(1)$ slack.

This proves all finite-dimensional path, compactly supported mark,
and slack limits jointly. Lemma~\ref{lem:bridge},
Proposition~\ref{prop:polypp}, and Proposition~\ref{prop:slack}
give tightness of the three marginals, hence of their product.
The resulting joint weak limit is exactly \eqref{eq:joint-main}.
The top-point continuity argument in Corollary~\ref{cor:orderstats}
proves \eqref{eq:argmax} and its stated independence properties.
\end{proof}

\section{Completion of the qualitative lacunary point-process question}
\label{sec:lacunary}

For lacunary weights the extreme locations are spread over a positive
fraction of the variance clock, unlike the polynomial case. The
negligible-prefix argument alone would therefore not prove the
independence theorem. Rare-mark factorization replaces it.

\begin{proposition}[Lacunary marks under the tilted product]
\label{prop:lacpp}
Under \eqref{eq:lacunarity},
\begin{equation}
 (\mathcal B_t,\Xi_t^{\rm lac})\Rightarrow(B,\Xi)
 \quad\text{under }Q_t, \label{eq:tiltedlacjoint}
\end{equation}
with independence. The same assertion holds when marks are restricted
to $j\le\lfloor cn\rfloor$, with limiting intensity restricted to
$0\le u\le c$, for fixed $0<c<1$.
\end{proposition}
\begin{proof}
Fix a height window $[L_0,L_1]$. For $j\le(1-\eps)n$,
$a_j\ge q^{-\eps n+O(1)}$, so its cap exceeds $\log n+L_1$
and its normalization $1-e^{-a_j}$ tends uniformly to one.
With location $u=j/n$ and height $Y_j-\log n=z$, the intensity
calculation on this prefix is therefore the Riemann sum
\[
 \frac1n\sum_{j\le(1-\eps)n}
    \int(1-e^{-g(j/n,z)})e^{-z}\,dz.
\]
Its limit is the required intensity integral over $[0,1-\eps]$.
For the remaining active indices, the probability of a contribution
in this height window is at most
\[
 \frac{e^{-L_0}}{n(1-e^{-\log n-L_0})}.
\]
Their combined contribution is at most $C_g\eps+O_g(n^{-1})$.
Inactive coordinates have cap below one and cannot contribute for
large $n$. Let $\eps\downarrow0$. The sum of the probabilities
has the Poisson intensity limit and their maximum is $O_g(n^{-1})$,
so the independent-product Laplace argument proves
$\Xi_t^{\rm lac}\Rightarrow\Xi$. Restricting to $j\le cn$ proves
the restricted statement.

Proposition~\ref{prop:lacdimension} identifies the index and variance
clocks up to $O_q(n^{-1})$. Lemma~\ref{lem:bridge} and
Lemma~\ref{lem:factorization}, now with $c_t=\log n$ and
$V\sim n$, prove independent Brownian and Poisson limits. The
moment bounds and the local count bounds give joint tightness.
\end{proof}

\begin{proof}[Proof of Theorem~\ref{thm:lacunary}]
First retain only marks with $j\le cn$, where $1/2<c<1$ is fixed.
For finitely many clock times strictly below $c$, observe the whole
prefix $I=\{1,\ldots,\lfloor cn\rfloor\}$. By
\eqref{eq:lacclock}, its variance fraction tends to $c$.
The tilted joint law of its Gaussian vector, its total, and its marks
is described by Proposition~\ref{prop:lacpp}. Applying
\eqref{eq:jointlikelihood} and Lemma~\ref{lem:llt} yields the
limiting weight $e^{-h}g_c(B(c))$. Thus the marks on $[0,c]$ remain
Poisson and independent, the Gaussian component becomes a bridge,
and the slack becomes independent exponential.

It remains to remove the restriction on the marks, not to assume it
harmless. Put $J=\{j>\lfloor cn\rfloor\}$. Its complement has at
least a fixed positive fraction of the total variance, uniformly for
$c\ge1/2$, so its $P_x$ marginal likelihood is bounded by a constant
independent of $c$. Formula \eqref{eq:exceedance} and the fact that
only the active indices contribute show, for each fixed $L_0$,
\begin{align}
 \limsup_{t\to\infty}P_x\big(
  \exists j>cn:\ Y_j-\log n>L_0\big)
 &\le C_q(1-c)e^{-L_0}. \label{eq:removeend}
\end{align}
Indeed the corresponding tilted union bound has at most
$(1-c)n+1$ terms, each bounded by
$e^{-L_0}/[n(1-e^{-\log n-L_0})]$; all inactive terms vanish.
The limiting Poisson process has the analogous vanishing probability
for a point in $(c,1]\times(L_0,\infty)$.

For any fixed finite set of path times below one, choose $c$ larger
than all of them, use the truncated joint result, and then let
$c\uparrow1$ in \eqref{eq:removeend}. This gives the full joint
finite-dimensional and compact-mark limits. The path and slack
components are tight by Lemma~\ref{lem:bridge}; the truncated-process
approximation and \eqref{eq:removeend} give tightness for the marks.
Consequently the triple converges to the mutually independent triple
in \eqref{eq:lacjoint}.

For completeness, the highest mark cannot escape to arbitrarily large
height under $P_x$. Divide the active indices into two half-blocks.
Each has a uniformly bounded marginal likelihood because its
complement has asymptotic variance fraction $1/2$. Under $Q_t$, the
expected number above $\log n+L$ is at most $C e^{-L}$, so the same
upper bound, with another constant, holds under $P_x$. All marks above
a fixed height have locations at most one. At low heights the limiting
Poisson count grows without bound. The continuity argument for the top
points therefore applies just as in Corollary~\ref{cor:orderstats},
giving the uniform location and independent Gumbel height.
\end{proof}

\subsection{What changes between the two weight regimes}
There is one universal limiting intensity but two different spatial
normalizations. In the lacunary case an interior active coordinate has
a cap much larger than $\log n$, and essentially all $n$ active
coordinates compete for the maximum. In the polynomial case the cap
at $j=Nu$ is $u^{-p}$, a quantity of order one. To compete at height
$b\sim\log N$, an index must satisfy
$j\lesssim N b^{-1/p}$. This moving cap boundary is why a naive
centering by $\log N$ fails and why the location scale is $M$, not $N$.

The common bridge arises from the global mass constraint. The Poisson
points survive independently because finitely many extreme contributions
are only logarithmic in size, negligible on the square-root variance
scale. Lemma~\ref{lem:factorization} turns that statement into a joint
characteristic-function calculation rather than an independence heuristic.

\section{Reproducible symbolic and numerical diagnostics}
\label{sec:numerics}

The archive contains \texttt{code/verify.py}, the generated table
\texttt{data/maximum\_table.csv}, and a machine-readable execution record
\texttt{data/verification.json}. The program performs six exact symbolic
checks: the identity $v=m-am'$ and the five coefficient identities in
\eqref{eq:allorder} for orders zero through four. All six passed in the
delivered run. It then evaluates the independent maximum distribution
exactly as a finite mathematical product, and evaluates the conditional
normalizers by numerical Fourier inversion with an analytically
represented infinite tail.

All decimal evaluations are floating-point diagnostics, not interval
certificates and not a substitute for the proofs. In particular, agreement
under changes of numerical settings does not establish a bound on all
rounding, truncation, or quadrature errors.

\subsection{The exact formulas evaluated}
For $y>0$, let $K_y=\lfloor Ny^{-\beta}\rfloor$. The independent
maximum formula is
\begin{equation}
 \log Q_t(\max_jY_j\le y)
 =K_y\log(1-e^{-y})
   -\sum_{j=1}^{K_y}\log(1-e^{-(N/j)^p}).
 \label{eq:numericalQ}
\end{equation}
An equality case $a_j=y$ contributes zero and causes no ambiguity.
No infinite truncation is needed in this formula.

If $\Phi_T(s)=\E_{Q_t}e^{is(T-\mu)}$, Fourier inversion and
\eqref{eq:D} give
\begin{equation}
 D=\frac1\pi\int_0^\infty
           \Re\frac{\Phi_T(s)}{1-is}\,ds. \label{eq:fourierD}
\end{equation}
The analogous $D_y$ formula uses the capped law but remains centered
at the original $\mu$. Its characteristic function is the centered
capped characteristic function multiplied by $e^{-is\delta_y}$.
The density smoothness established in Lemma~\ref{lem:llt} justifies
the inversion and the exponentially integrable slack kernel.

The code evaluates the first $J=\lceil4N\rceil$ coordinates directly.
For the rest, it uses the convergent expansion
\begin{equation}
 \ell(a):=\log\frac{1-e^{-a}}a
 =-\frac a2+\sum_{\substack{r\ge2\\r\ \mathrm{even}}}
       c_ra^r,\qquad c_r=\frac{B_r}{r\,r!},\quad |a|<2\pi.
 \label{eq:bernoullitail}
\end{equation}
Here $B_r$ are the Bernoulli numbers. The centered characteristic
logarithm for a tail coordinate is
\[
 \sum_{\substack{r\ge2\\r\ \mathrm{even}}}
 c_ra^r\big((1-is)^r-1+irs\big).
\]
Summation over $j>J$ uses the exact Hurwitz-zeta identity
\begin{equation}
 \sum_{j>J}a_j^r=N^{pr}\zeta(pr,J+1). \label{eq:zetatail}
\end{equation}
The linear term cancels on centering; the tail mean and variance are
calculated from the same expansion. The delivered run retains the even
terms through order $16$ and integrates in the scaled variable
$u=s\sqrt V$ up to $u=30$. It does not represent discarded coordinates
by zero, which would introduce an uncontrolled shift on the small-sum
scale.

\subsection{Results and finite-size effects}
Table~\ref{tab:central} reports $p=2$, $z=0$. The conditional column
uses \eqref{eq:capidentity} and \eqref{eq:fourierD}; the independent
column uses \eqref{eq:numericalQ}. The incomplete-gamma column is
$\exp[-N\Gamma(1/2,b)]$. Their common limiting Gumbel value is
$e^{-1}=0.3678794412\ldots$.

\begin{table}[htbp]
\centering
\caption{Central maximum probabilities for polynomial weights $j^{-2}$.}
\label{tab:central}
\begin{tabular}{rrrrr}
\toprule
$N$ & $b$ & Conditional $P_x$ & Tilted $Q_t$ & Gamma approximation\\
\midrule
32   & 2.928492 & 0.450850950 & 0.415315863 & 0.414779547\\
128  & 4.141501 & 0.413874108 & 0.403560486 & 0.403369337\\
512  & 5.395538 & 0.399337502 & 0.396342140 & 0.396309268\\
2048 & 6.675404 & 0.392427364 & 0.391548647 & 0.391540557\\
\bottomrule
\end{tabular}
\end{table}

The full CSV additionally contains $z=-1,1$ at all four dimensions,
for twelve parameter triples. At $N=2048$ and $z=-1$, for example,
the conditional probability is approximately $0.060955084$, the
independent probability $0.064524504$, and the gamma approximation
$0.064728174$. The conditioning-to-tilt ratio is about $0.944681$.
Thus the table does not conceal a sizable lower-tail finite-size effect
behind the limiting theorem. The constants implicit in uniform bounds
may depend on the chosen bounded $z$-window.

An independent stability rerun at $p=2,N=128,z=0$ changed $J$ from
$4N$ to $6N$, the tail order from $16$ to $20$, and the scaled
integration cutoff from $30$ to $36$. The absolute change in the
conditional CDF was $5.0\times10^{-16}$ in the delivered environment.
The mean and variance changed by approximately
$2.8\times10^{-14}$ and $1.4\times10^{-14}$, respectively.
An independent high-precision integral evaluation of $A_2$ agreed
with the gamma--zeta expression to about $2.6\times10^{-25}$.
These are observations about numerical stability, not rigorous error
bars on Table~\ref{tab:central}.

The executed environment was Python~3.13.5, NumPy~2.3.5,
SciPy~1.17.0, SymPy~1.14.0, and mpmath~1.3.0. The script records the
principal versions in its JSON output. Reproduction does not require
GitHub access or downloading the repository.

\clearpage
\section{Proof status, formalization, and limitations}
\label{sec:status}

\subsection{A precise status ledger}
\begin{longtable}{>{\raggedright\arraybackslash}p{.34\linewidth}
                  >{\raggedright\arraybackslash}p{.59\linewidth}}
\toprule
Component & Status and scope\\
\midrule
\endhead
Uniform-series tilting and Gibbs conditioning & Classical framework;
needed density and likelihood calculations rederived here.
The general variance-fraction principle is already in \cite{repounif}.\\[4pt]
Polynomial saddle constants and variance clock & Established background
for the extension, proved here for self-containment; not advertised as
new relative to the repository.\\[4pt]
Incomplete-gamma conditional maximum law & Proved in
Theorem~\ref{thm:gamma}; quantitative error for fixed $p>1$ and bounded
centered height.\\[4pt]
Every fixed-order logarithmic correction & Proved with explicit
coefficient polynomials and a remainder in \eqref{eq:allorder}; no
growing-order uniformity claimed.\\[4pt]
Polynomial marked extremes and localization & Proved in
Theorem~\ref{thm:jointpoly}; spatial scale $N/b^{1/p}$ and all fixed
upper order statistics.\\[4pt]
Joint independence of bridge, points, and slack & Proved for polynomial
and uniformly lacunary weights, with the displayed process topologies.\\[4pt]
Lacunary Question~7 & Its qualitative point-process and independence
parts are answered. Its explicit joint-rate request is not answered.\\[4pt]
Numerical work & Six symbolic identities passed; twelve floating-point
CDF diagnostics and a stability rerun were executed. No interval or
formal analytic certification of the decimal outputs.\\[4pt]
Lean/Rocq verification & Not performed; no compiled formal theorem or
repository build is claimed or included.\\[4pt]
Global originality and peer review & Not established. Candidate
contributions are identified relative to the inspected sources, not
certified as historical firsts.\\
\bottomrule
\end{longtable}

\subsection{A realistic formalization path}
The most economical first target is the maximum theorem rather than
the complete point-process topology. The finite-product identity
\eqref{eq:numericalQ}, the integral comparison
\eqref{eq:lattice}, and the exact gamma integral
\eqref{eq:gammaint} form a small, independently useful layer. These
statements involve finite sums, a monotone real integral, and a
single change of variables; they do not depend on the unformalized
conditioning manuscripts.

The next layer should formalize the tilted truncated-exponential mean
and variance, the product-capping identity \eqref{eq:capidentity},
and a reusable density local limit theorem with explicit hypotheses.
The proof boundary must state whether Fourier inversion, the normal
characteristic function, and the moment tightness criterion are imported
library theorems or independently proved. Merely defining an abstract
density that is assumed to obey \eqref{eq:llt} would make the final
maximum theorem conditional on that assumption.

The bridge and Poisson layer then requires probability kernels on a
countable product, the likelihood disintegration
\eqref{eq:jointlikelihood}, weak convergence of continuous paths,
and Laplace-functional convergence for counting measures. The rare-mark
factorization lemma is a useful finite-product starting point because
its core estimate is a direct telescoping inequality. Numerical output
should remain outside the trusted proof boundary unless a certified
checker is supplied for the precise numerical assertion.

\subsection{Limits of the statements}
The constants in the polynomial estimates are not asserted uniform as
$p\downarrow1$ or $p\to\infty$. The weights are exactly $j^{-p}$
for the gamma expansion; a slowly varying multiplier can alter its
correction terms. The joint process convergence is qualitative, even
though the scalar maximum law is quantitative. The point-process
assertions use vague topology, not total variation between a discrete
location lattice and a continuously located Poisson process; the latter
comparison would be singular without coarsening or randomizing locations.

The proofs address the inequality $S\le x$. A regular conditional law
on $S=x$, a thin shell, or several simultaneous constraints requires
its own likelihood and normalization. No such extension is silently
included. Finally, the manuscript has not been externally refereed;
the distinction between a complete argument written here and an
independently checked theorem remains important.

\section{Further research questions}
\label{sec:questions}

\paragraph{1. Quantitative joint convergence.}
Complete the remaining rate component of the lacunary report's
Question~7. Specify a metric combining path approximation with
coarsened or bounded-Lipschitz point-process observables, and derive
an explicit bound for joint independence. Exact total variation of
continuous locations versus the deterministic index lattice is not an
appropriate target. The rare-mark factorization proof already exposes
an $O(\log R/\sqrt R)$ cross-term before conditioning; controlling the
complement likelihood and the final clock segment is the next obstacle.

\paragraph{2. The first genuine conditioning correction.}
Replace the conservative $O(N^{-1/2})$ error in
Theorem~\ref{thm:gamma} by an Edgeworth calculation for $D_y/D$,
uniformly on a bounded height window. The exact quantities
$\delta_y=\mu-\mu_y$ and $V-V_y$ identify the terms that must enter.
Can one prove an explicit leading correction at order
$(\log N)^2/N$, with a smaller controlled remainder? This is separate
from the inverse-logarithmic series already proved.

\paragraph{3. Regularly varying weights.}
For decreasing $w_j=j^{-p}\ell(j)$ with slowly varying $\ell$, seek
a marked-extreme theorem in terms of the counting function
$\nu_t(y)=\#\{j:tw_j>y\}$. The natural first centering equation is
$e^{-b_t}\nu_t(b_t)\approx1$. Determine the second-order regular
variation assumptions that recover explicit corrections and distinguish
them from the special exact gamma structure of pure powers.

\paragraph{4. Transition between polynomial and lacunary geometry.}
Study arrays in which the decay exponent or local log-slope changes
with $t$. The limits $p\to\infty$ and $N\to\infty$ need not commute.
Find a uniform description interpolating between localization in
$o(N)$ indices and competition across a positive fraction of all
active coordinates. The endpoint $p\downarrow1$ has its own
summability and scale issues and should be analyzed separately.

\paragraph{5. Nonuniform microscopic laws.}
Replace the uniforms by variables with density asymptotic to
$c u^{\kappa-1}$ at zero. Their strongly tilted coordinates should be
compared with gamma, rather than exponential, variables. Determine
which regularity at zero and which global upper-tail assumptions
suffice for joint Poisson--bridge independence. Even the logarithmic
centering must account simultaneously for the gamma tail and the
moving cap boundary; it should not be copied from \eqref{eq:scales}.

\paragraph{6. Thin shells and multiple constraints.}
For $x-h_t\le S\le x$, determine how the limit depends on $th_t$,
including a precise version of conditioning at $S=x$. With several
weighted-sum constraints, formulate nondegeneracy hypotheses under
which the Gaussian limit is projected off a finite-dimensional
constraint space. Identify when the extremes remain independent and
when a constraint directly detects their locations.

\paragraph{7. Random weights and geometric phase.}
Establish quenched and annealed counterparts under random lacunarity
or random regular variation. For geometric weights, combine the
marked extremes with the phase-dependent microscopic boundary law
already in \cite{repolac,repounif}. Joint phase mixing must be proved
at the level of distributions; averaging a parameter inside a density
is not generally the same as averaging the density itself.

\paragraph{8. Certified analytic computation.}
Turn \eqref{eq:fourierD}--\eqref{eq:zetatail} into an interval-certified
algorithm with explicit Fourier-tail and Bernoulli-tail bounds,
then determine bit complexity for a requested absolute or relative
CDF error. The current numerical stability checks do not provide
that certificate. A formal finite-product and gamma-integral layer
would be a natural first contribution back to the repository.

\section{Conclusion}
The rare-event constraint has two mathematically distinct effects.
It produces a Brownian bridge from collective square-root fluctuations
and an exponential boundary slack, but it leaves logarithmic extremes
asymptotically Poisson and independent of those objects. The locations
of those extremes retain detailed information about the weight decay.
They fill the active index range for uniformly lacunary weights and
localize on $N/(\log N)^{1/p}$ indices for polynomial weights.

For pure polynomial weights, the entire fixed-order logarithmic
maximum expansion is encoded by a single upper incomplete gamma
integral. The key quantitative step is not the gamma asymptotics by
itself, but the estimate showing that reimposing the small-sum
constraint changes the logarithmic maximum probability by only
$O(N^{-1/2})$ on bounded centered windows. The lacunary joint theorem
answers a specifically identified qualitative research question in
the inspected repository; the polynomial gamma law and localization
provide a separate, more quantitative extension.

\appendix
\section{Reproduction and source provenance}
\label{app:reproduce}
The archive is self-contained for TeX compilation. It contains the
source \texttt{article.tex}, its compiled PDF, a \texttt{Makefile},
\texttt{README.md}, \texttt{CLAIM\_STATUS.md}, the verification script,
a requirements file, and generated numerical data. A standard TeX Live
installation with the AMS packages, \texttt{newtx}, \texttt{microtype},
and \texttt{xurl} is sufficient. Run from the archive directory:
\begin{verbatim}
make pdf
python -m pip install -r requirements.txt
python code/verify.py
\end{verbatim}
The optional \texttt{--quick} flag omits $N=2048$ and the stability
rerun. The full run regenerates the JSON and CSV data. It does not
fetch external mathematical sources and does not modify the repository.

The repository comparison was made at \commit. The two pinned
manuscripts and the newer manuscript's claim-status file were
inspected; their earlier snapshot identifiers are provenance inside
those manuscripts, not the snapshot used here. Narrow repository
searches were used to locate related work. This is not an assertion
that every file in the repository has been read or that all possible
overlap has been excluded. The external comparison used the primary
publisher records for \cite{df,dz}, Arias de Reyna's paper, and the
NIST special-function reference. The maximum, point-process, and
joint-independence proofs above are supplied explicitly rather than
inferred from those bibliographic records.

\clearpage
\begin{thebibliography}{9}
\bibitem{arias}
J.~Arias de Reyna,
\emph{An infinitely differentiable function with compact support:
Definition and properties}.
English translation of the original article in
Revista de la Real Academia de Ciencias Exactas, F\'isicas y Naturales
\textbf{76} (1982), 21--38.
\href{https://arxiv.org/abs/1702.05442}{arXiv:1702.05442} (2017).
Theorem~3 gives the uniform random-series interpretation.

\bibitem{df}
P.~Diaconis and D.~A.~Freedman,
\emph{Conditional limit theorems for exponential families and finite
versions of de Finetti's theorem}.
Journal of Theoretical Probability \textbf{1} (1988), 381--410.
\href{https://doi.org/10.1007/BF01048727}{doi:10.1007/BF01048727}.

\bibitem{dz}
A.~Dembo and O.~Zeitouni,
\emph{Refinements of the Gibbs conditioning principle}.
Probability Theory and Related Fields \textbf{104} (1996), 1--14.
\href{https://doi.org/10.1007/BF01303799}{doi:10.1007/BF01303799}.

\bibitem{dlmf}
NIST Digital Library of Mathematical Functions,
Chapter~8, \emph{Incomplete Gamma and Related Functions}, and
Chapter~25, \emph{Zeta and Related Functions}.
\href{https://dlmf.nist.gov/8}{Chapter 8};
\href{https://dlmf.nist.gov/25.5}{Section 25.5, integral representations}.
Consulted 30 September 2026. The specific gamma expansion used in this
article is rederived in \eqref{eq:gamma-shift}.

\bibitem{repolac}
\repo\ research corpus,
\emph{Sharp Conditioning Laws for Fabius and Lacunary Random Series:
A global simplex approximation, exact total-variation thresholds,
and phase-dependent boundary layers}.
Research manuscript prepared for V.~Reshetnikov, 28 September 2026.
Inspected at commit \commit.
\href{https://github.com/VladimirReshetnikov/ProveIt/blob/a4268e78ebd0f6bf71ba609c07f4b3415748f532/Analysis/FabiusFunction/docs/semi-formalized-research-frontiers/drafts/representations/Sharp_Conditioning_Laws_Fabius_Lacunary_Series/article.tex}{Pinned manuscript source}.
Question~7 is the qualitative point-process and independence target.

\bibitem{repounif}
\repo\ research corpus,
\emph{Sharp Conditioning Laws for Uniform Random Series:
Variance thresholds, boundary phases, Brownian bridges, and rare-event
sampling}.
Research manuscript prepared for V.~Reshetnikov, 28 September 2026.
Inspected at commit \commit.
\href{https://github.com/VladimirReshetnikov/ProveIt/blob/a4268e78ebd0f6bf71ba609c07f4b3415748f532/Analysis/FabiusFunction/docs/semi-formalized-research-frontiers/drafts/representations/Sharp_Conditioning_Laws_Uniform_Random_Series/article.tex}{Pinned manuscript source};
\href{https://github.com/VladimirReshetnikov/ProveIt/blob/a4268e78ebd0f6bf71ba609c07f4b3415748f532/Analysis/FabiusFunction/docs/semi-formalized-research-frontiers/drafts/representations/Sharp_Conditioning_Laws_Uniform_Random_Series/CLAIM_STATUS.md}{pinned claim-status ledger}.
The general variance criterion and polynomial clocks are existing
repository results, not new claims of this article.
\end{thebibliography}
\end{document}
